The 14 states that use independent or bipartisan redistricting commissions represent the most politically established infrastructure for neutral redistricting in the United States. California, Colorado, Michigan, Arizona, Hawaii, Idaho, Montana, New Jersey, New York, Ohio, Virginia, Washington, and Alaska all use commissions for at least one chamber; Iowa uses a non-commission variant (the Legislative Services Agency model) that achieves similar neutrality through different institutional design.

Commission redistricting has significant advantages over legislative redistricting: commissions are structurally independent, typically required to hold public hearings, and must satisfy enumerated criteria rather than partisan objectives. But commissions are not gerrymander-proof. Even well-designed commissions can produce maps that favor one party through the cumulative effect of individually defensible line-drawing decisions. The California CCRC's 2021 congressional maps, for example, were criticized by some analysts for producing more Democratic-leaning seats than a neutral algorithm would generate---not because the commission acted with partisan intent, but because the commission's good-faith application of criteria (particularly community of interest) has distributional consequences that are not transparent to the commission.

The \texttt{bisect} platform addresses this gap. By providing a deterministic baseline that any person can verify, it makes the distributional consequences of commission criteria explicit: the difference between the algorithmic baseline and the commission's final map is the quantifiable effect of the commission's discretionary choices. This transparency serves two goals. It holds commissions accountable by making manipulation visible. And it protects commissions from unfair criticism by proving that departures from neutrality are VRA-required or otherwise legally justified.

This paper is addressed to redistricting commissioners, state administrators, and commission legal counsel. Section~2 examines the Iowa model as the strongest existing analogue. Section~3 describes three integration modes. Section~4 provides a detailed Commission Adjustment Protocol for Mode~(b) integration. Section~5 concludes.
